\section{Proposed JSON Standard for MicroSim Performance Data}

This section proposes a concrete JSON schema for storing MicroSim performance
data. The design has three goals: to capture the fields necessary to compute
the fitness function of Section~\ref{sec:fitness}; to align with the
Experience API so that the same instrumentation works in both deployment
phases; and to support the archival, variant-centric data model that the
evolutionary loop requires.

\subsection{Design Principles}

The schema follows several principles drawn from the xAPI and
cmi5~\cite{aicc2015cmi5} specifications and from the requirements of the
evolutionary loop:

\begin{itemize}
    \item \textbf{Variant-centric.} Every record identifies the specific
          MicroSim variant, not merely the concept, so that fitness can be
          computed per variant. This is the field that makes evolution possible
          and is the most important addition over conventional analytics.
    \item \textbf{xAPI-aligned.} Field names and structure mirror the xAPI
          Actor--Verb--Object--Result pattern. The xAPI specification recommends
          using built-in properties (such as the Result score) before resorting
          to custom extensions, since custom extensions may be excluded from
          standard LRS reports; the schema follows this guidance.
    \item \textbf{Standard verbs.} The schema uses registered xAPI and cmi5
          verbs such as \texttt{initialized}, \texttt{interacted},
          \texttt{answered}, \texttt{completed}, \texttt{passed}, and
          \texttt{failed} rather than inventing new ones, preserving
          interoperability.
    \item \textbf{Privacy-conscious.} Learner identity is carried in an
          account-based identifier rather than personal information, and the
          schema supports a fully anonymous mode for Phase~1.
\end{itemize}

\subsection{The MicroSim Event Record Schema}

The following is the proposed per-event record. It can be emitted directly
to a web analytics platform in Phase~1 (flattened) or wrapped in a full xAPI
statement in Phase~2.

\begin{lstlisting}[style=json]
{
  "$schema": "https://intelligenttextbooks.org/schemas/microsim-event/v1.json",
  "event_id": "a7f3c918-b88b-4220-90a5-a4c32391a240",
  "schema_version": "1.0.0",
  "timestamp": "2026-06-14T15:42:08.000Z",

  "actor": {
    "mode": "identified",
    "account_id": "student-uuid-or-anonymous-session-hash",
    "segment": "ib-physics-sl"
  },

  "verb": {
    "id": "http://adlnet.gov/expapi/verbs/completed",
    "display": "completed"
  },

  "microsim": {
    "sim_id": "bouncing-ball",
    "variant_id": "bouncing-ball-v3-graph-overlay",
    "parent_variant_id": "bouncing-ball-v2",
    "generation": 3,
    "hypothesis": "Real-time velocity graph improves transfer to symbolic problems",
    "source_url": "https://example.github.io/sims/bouncing-ball-v3/main.html"
  },

  "concept": {
    "concept_id": "newtons-second-law",
    "graph_node": "phys.mechanics.dynamics.n2l",
    "bloom_level": 2,
    "bloom_label": "Apply"
  },

  "result": {
    "score": {
      "scaled": 0.75,
      "raw": 15,
      "min": 0,
      "max": 20
    },
    "success": true,
    "completion": true,
    "duration_iso8601": "PT2M22S",
    "duration_seconds": 142
  },

  "interaction": {
    "attempts": 2,
    "param_changes": 11,
    "controls_used": ["drop_height", "gravity"],
    "idle_seconds": 8,
    "engagement_depth": 0.82
  },

  "learning": {
    "pre_assessment_score": 0.40,
    "post_assessment_score": 0.75,
    "normalized_gain": 0.58,
    "prediction_made": true,
    "prediction_correct": false
  },

  "context": {
    "textbook_id": "ib-physics-sl",
    "chapter": "dynamics",
    "registration": "ec231277-b27b-4c15-8291-d29225b2b8f7",
    "platform": "mkdocs-material",
    "client": "web",
    "ab_test_group": "treatment-graph-overlay"
  }
}
\end{lstlisting}

\subsection{Field Reference}

The schema's blocks correspond to the components of the fitness function and
to xAPI's statement structure:

\begin{itemize}
    \item \textbf{\texttt{actor}} carries learner identity. The \texttt{mode}
          field switches between \texttt{anonymous} (Phase~1) and
          \texttt{identified} (Phase~2). The \texttt{segment} enables
          audience-conditioned evolution as described in
          Section~\ref{sec:compound-interest}.
    \item \textbf{\texttt{verb}} uses a registered xAPI verb identifier and
          display string, ensuring LRS interoperability.
    \item \textbf{\texttt{microsim}} is the variant-centric core. The
          \texttt{variant\_id}, \texttt{parent\_variant\_id}, and
          \texttt{generation} fields encode the evolutionary lineage directly,
          allowing the archive to be reconstructed as a tree. The
          \texttt{hypothesis} field documents \textit{why} the variant was
          created, paralleling the DGM's practice of recording what has been
          tried and why.
    \item \textbf{\texttt{concept}} ties the event to the dependency graph and
          Bloom's level, supplying the structured search-space coordinates.
    \item \textbf{\texttt{result}} uses xAPI's standard \texttt{score} object
          (scaled, raw, min, max), \texttt{success}, \texttt{completion}, and
          ISO~8601 duration, in keeping with the recommendation to prefer
          built-in properties.
    \item \textbf{\texttt{interaction}} captures engagement-quality signals
          that distinguish meaningful exploration from disengagement.
    \item \textbf{\texttt{learning}} carries the pre/post assessment and
          normalized-gain fields that drive the primary fitness component.
    \item \textbf{\texttt{context}} supplies textbook, chapter, A/B group, and
          the cmi5-style \texttt{registration} identifier that groups statements
          within a learning session.
\end{itemize}

\subsection{The Aggregated Variant Fitness Record}

Individual events are periodically aggregated into a per-variant fitness
record. This is the structure the selection step reads from, and it serves as
the sidecar metadata stored alongside each MicroSim variant in the archive.

\begin{lstlisting}[style=json]
{
  "$schema": "https://intelligenttextbooks.org/schemas/variant-fitness/v1.json",
  "variant_id": "bouncing-ball-v3-graph-overlay",
  "concept_id": "newtons-second-law",
  "generation": 3,
  "parent_variant_id": "bouncing-ball-v2",
  "hypothesis": "Real-time velocity graph improves transfer to symbolic problems",
  "window": {
    "start": "2026-05-01T00:00:00Z",
    "end": "2026-06-01T00:00:00Z"
  },
  "sample": {
    "n_learners": 412,
    "n_events": 1187,
    "segments": {
      "ib-physics-sl": 318,
      "us-high-school": 94
    }
  },
  "metrics": {
    "avg_normalized_gain": 0.58,
    "avg_post_score": 0.75,
    "bloom_advancement_rate": 0.41,
    "avg_attempts": 2.1,
    "avg_time_seconds": 138,
    "avg_engagement_depth": 0.79,
    "retention_score": 0.68,
    "completion_rate": 0.91
  },
  "fitness": {
    "composite_score": 0.712,
    "weights": {
      "normalized_gain": 0.35,
      "bloom_advancement": 0.20,
      "efficiency": 0.10,
      "engagement_quality": 0.10,
      "retention": 0.25
    },
    "confidence_interval_95": [0.681, 0.743],
    "rank_within_concept": 1
  },
  "status": "active",
  "human_review": {
    "approved": true,
    "reviewer": "content-team",
    "reviewed_at": "2026-05-02T10:15:00Z"
  }
}
\end{lstlisting}

The \texttt{fitness} block records both the composite score and the weights
used to compute it, so that scoring remains auditable and reproducible as
priorities change. The \texttt{confidence\_interval\_95} field is essential:
selection should account for statistical uncertainty, preferring variants whose
advantage is well-established over those with a high but noisy point estimate.
The \texttt{human\_review} block enforces the oversight requirement discussed
in Section~\ref{sec:safety}; a variant cannot reach \texttt{active} status
without review.

\subsection{The Full xAPI Statement Form (Phase~2)}

For Phase~2, the event record maps onto a conformant xAPI statement. The
MicroSim-specific fields move into result and context extensions, with custom
extension identifiers under an organizational namespace, while standard fields
use built-in properties~\cite{rustici2024cmi5}:

\begin{lstlisting}[style=json]
{
  "id": "a7f3c918-b88b-4220-90a5-a4c32391a240",
  "actor": {
    "objectType": "Agent",
    "account": {
      "homePage": "https://example.github.io",
      "name": "student-uuid"
    }
  },
  "verb": {
    "id": "http://adlnet.gov/expapi/verbs/completed",
    "display": { "en-US": "completed" }
  },
  "object": {
    "id": "https://example.github.io/sims/bouncing-ball-v3",
    "objectType": "Activity",
    "definition": {
      "name": { "en-US": "Bouncing Ball MicroSim v3 (graph overlay)" },
      "type": "http://adlnet.gov/expapi/activities/simulation",
      "extensions": {
        "https://intelligenttextbooks.org/xapi/ext/variant-id":
            "bouncing-ball-v3-graph-overlay",
        "https://intelligenttextbooks.org/xapi/ext/concept-id":
            "newtons-second-law",
        "https://intelligenttextbooks.org/xapi/ext/bloom-level": 2
      }
    }
  },
  "result": {
    "score": { "scaled": 0.75, "raw": 15, "min": 0, "max": 20 },
    "success": true,
    "completion": true,
    "duration": "PT2M22S",
    "extensions": {
      "https://intelligenttextbooks.org/xapi/ext/normalized-gain": 0.58,
      "https://intelligenttextbooks.org/xapi/ext/attempts": 2,
      "https://intelligenttextbooks.org/xapi/ext/engagement-depth": 0.82
    }
  },
  "context": {
    "registration": "ec231277-b27b-4c15-8291-d29225b2b8f7",
    "contextActivities": {
      "category": [
        { "id": "https://w3id.org/xapi/cmi5/context/categories/moveon" }
      ],
      "parent": [
        { "id": "https://example.github.io/textbooks/ib-physics-sl/dynamics" }
      ]
    }
  },
  "timestamp": "2026-06-14T15:42:08.000Z"
}
\end{lstlisting}

This statement is conformant with the xAPI data model and follows cmi5
conventions for category context activities and the registration identifier,
so it can be stored in any standards-compliant Learning Record Store and
queried by the fitness-aggregation pipeline.
